\chapter{Wie wirkt Magie?}

\section{Was ist Magie?}

Magie gehört zu den außergewöhnlichsten Fähigkeiten, die ein Held in Heldenpfade erlernen kann. Während viele Helden ihre Herausforderungen mit Mut, Geschick oder körperlicher Stärke überwinden, greifen Magiebegabte auf geheimnisvolle Kräfte zurück, um die Welt um sich herum zu beeinflussen.

Magie kann Verletzungen heilen, Feuer beschwören, Verbündete stärken oder Gegner schwächen. Manche Zauber wirken sofort, andere entfalten ihre Wirkung über längere Zeit oder müssen aufrechterhalten werden.

Nicht jeder Held ist magiebegabt. Die meisten Helden verlassen sich ausschließlich auf ihre Waffen, ihre Ausrüstung und ihre Talente. Erst das Talent \emph{Magiebegabt} ermöglicht einem Helden überhaupt den Zugang zur Magie.

Die einzelnen Zauber unterscheiden sich stark in ihrer Wirkung. Die allgemeinen Regeln zum Wirken von Magie werden in diesem Kapitel beschrieben. Die Beschreibungen aller verfügbaren Zauber befinden sich in Anhang 2 -- Zauber.

\section{Mana}

Jeder Zauber kostet Mana. Mana beschreibt die innere magische Kraft eines Helden und bestimmt, wie viele Zauber er wirken kann, bevor er sich ausruhen muss.

Jeder Held besitzt zwei Werte:

\begin{itemize}
  \item \textbf{Manavorrat:} die maximale Menge an Mana, über die der Held verfügen kann.
  \item \textbf{Aktuelles Mana:} die Menge an Mana, die dem Helden momentan zur Verfügung steht.
\end{itemize}

Zu Beginn jedes Abenteuers besitzt ein Held stets seinen vollständigen Manavorrat als aktuelles Mana.

Immer wenn ein Held einen Zauber wirkt, bezahlt er dessen Manakosten. Das aktuelle Mana wird entsprechend reduziert.

Das aktuelle Mana kann niemals unter 0 sinken und niemals höher sein als der Manavorrat.

Verbrauchtes Mana wird normalerweise während einer Rast regeneriert. Die genauen Regeln hierzu befinden sich in Kapitel 8 -- Abenteuer erleben.

\section{Magiebegabt}

Um Zauber wirken zu können, muss ein Held das allgemeine Talent \emph{Magiebegabt} erlernt haben.

Ein Held ohne dieses Talent kann keine Zauber erlernen und keine Zauber wirken.

Das Talent verleiht dem Helden einen zusätzlichen Manavorrat von 6 Mana. Dadurch können bereits zu Beginn einer Kampagne mehrere Zauber erlernt werden.

Das Talent selbst verleiht noch keine Zauber. Welche Zauber ein Held beherrscht, entscheidet der Spieler beim Erlernen der Zauber.

\section{Zauber}

Jeder Zauber wird durch eine Zauberkarte beschrieben.

Eine Zauberkarte enthält alle Informationen, die zum Wirken des Zaubers benötigt werden. Dazu gehören insbesondere:

\begin{itemize}
  \item Name des Zaubers
  \item Manakosten
  \item Schlüsselwörter
  \item Reichweite
  \item Ziel
  \item Wirkungsdauer
  \item Beschreibung der Wirkung
\end{itemize}

Jeder Zauber besitzt eigene Eigenschaften. Manche Zauber wirken augenblicklich, andere müssen über längere Zeit aufrechterhalten werden oder können nur unter bestimmten Bedingungen eingesetzt werden.

Die vollständigen Beschreibungen aller Zauber befinden sich in Anhang 2 -- Zauber.
